% 张三慧《大学物理学（第4版）》V2，第16章《恒定电流》
% 原教材扫描PDF物理页118--120（印刷101--103页）。
% 根据原书数字及条件重述，保留16.6、16.11--16.14配图依赖。
\section*{第16章\quad 恒定电流}
\subsection*{章末习题16.1--16.18}

\paragraph{习题 16.1}
北京正负电子对撞机的电子储存环近似为
周长 \(240\,\mathrm m\) 的圆轨道，束流电流为
\(8\,\mathrm{mA}\)。电子速度近似光速，
估算环中同时运行的电子总数。

\paragraph{习题 16.2}
范德格拉夫静电加速器中的绝缘传送带宽
\(30\,\mathrm{cm}\)，速度 \(20\,\mathrm{cm/s}\)；
带面电荷密度足以在带的每一侧产生
\(1.2\times10^6\,\mathrm{V/m}\) 的电场。
求传送带输运的电流（用毫安表示）。
\textit{模型风险：}绝缘薄带一侧场强为
\(\sigma/(2\varepsilon_0)\)，不是导体表面的
\(\sigma/\varepsilon_0\)；未另行指定双面输运。

\paragraph{习题 16.3}
估算直径 \(1.0\,\mathrm{mm}\) 的银丝通过
\(30\,\mathrm A\) 电流时的电子漂移速度。
假定导电电子数与银原子数相等；
自己估算银的密度、摩尔质量等数量级参数。
并比较银丝 \(20^\circ\mathrm C\) 下
自由电子经典热运动平均速率。

\paragraph{习题 16.4}
一铜棒横截面
\(20\,\mathrm{mm}\times80\,\mathrm{mm}\)，
长 \(2\,\mathrm m\)，两端电势差
\(50\,\mathrm{mV}\)。
铜电阻率
\(\rho=1.75\times10^{-8}\,\Omega\mathrm m\)，
自由电子数密度
\(n=8.5\times10^{28}\,\mathrm{m^{-3}}\)。
求：（1）电阻；（2）电流；
（3）电流密度；（4）电场强度；
（5）耗散功率；（6）电子漂移速度。

\paragraph{习题 16.5}
一根铁制水管外直径 \(2.0\,\mathrm{cm}\)，
内直径 \(1.5\,\mathrm{cm}\)，输送自来水。
若沿管长方向流过总电流 \(20\,\mathrm A\)，
分别有多大电流经过铁管与管中的水？
水电阻率 \(0.01\,\Omega\mathrm m\)，
铁电阻率 \(8.7\times10^{-8}\,\Omega\mathrm m\)。

\paragraph{习题 16.6}
地下电话线由两根等长、均匀导线组成，
全长 \(L=5\,\mathrm m\)，两线在
图16.19的某点 \(P\) 短路。
测左端两端点 \(AB\) 的总电阻为
\(30\,\Omega\)，测右端 \(CD\) 电阻为
\(70\,\Omega\)，求故障点距左端的距离。
假定两端测量时没有其他并联漏电路径。

\paragraph{习题 16.7}
大气的少量正负离子使它具有弱导电性。
（1）晴朗地表附近平均电场
\(120\,\mathrm{V/m}\)，平均电流密度
\(4\times10^{-12}\,\mathrm{A/m^2}\)，
求等效大气电阻率。
（2）电离层与地面电势差
\(4\times10^5\,\mathrm V\)，
问对这两层间整体的总电阻有多大？
\textit{条件说明：}为得整体电阻需约定地表
面积或全球分布模型，下面按均匀穿过地球表面估算。

\paragraph{习题 16.8}
电缆由内半径
\(r_1=0.5\,\mathrm{cm}\) 的铜芯线和
外半径 \(r_2=2\,\mathrm{cm}\) 的同轴绝缘层组成，
绝缘层电阻率
\(\rho=10^{12}\,\Omega\mathrm m\)，
外面是金属铅护层（图16.20）。
求长 \(L=1000\,\mathrm m\) 电缆
沿径向的漏电电阻，以及铜芯与护层的电势差
\(100\,\mathrm V\) 时径向漏电电流。

\paragraph{习题 16.9}
同心球形电容器内外球半径分别为
\(r_1,r_2\)，夹层介质电阻率为 \(\rho\)。
证明漏电电阻
\[
 R=\frac{\rho}{4\pi}
 \left(\frac1{r_1}-\frac1{r_2}\right).
\]

\paragraph{习题 16.10}
一根输电线被风吹断后落地，
有 \(200\,\mathrm A\) 电流从接地点进入
近似平坦、均匀、各向同性的土壤，
土壤电阻率 \(10.0\,\Omega\mathrm m\)。
行人左脚距接地点 \(1.0\,\mathrm m\)，
右脚距其 \(1.3\,\mathrm m\)。
计算两脚所跨的地表电势差。
\textit{安全说明：}仅是理想半无限导体模型，
不是可依此接近高压事故现场的安全判断。

\paragraph{习题 16.11}
图16.21电路有电动势
\(\mathcal E_1=3.0\,\mathrm V\)、
内阻 \(r_1=0.5\,\Omega\)，
\(\mathcal E_2=6.0\,\mathrm V\)、
内阻 \(r_2=1.0\,\Omega\)，
及外部电阻 \(R_1=2.0\,\Omega\)、
\(R_2=4.0\,\Omega\)。
求通过 \(R_1,R_2\) 的电流。
\textit{原图连接：}左支路 \(R_1\)，中支路
\((\mathcal E_1,r_1)\)，
右支路上方串 \(R_2\)、
其下串 \((\mathcal E_2,r_2)\)，
两个电源长板都画在下侧；正负方向须按原图校核。

\paragraph{习题 16.12}
图16.22所示网络含
\(\mathcal E_1=3.0\,\mathrm V\)、
\(\mathcal E_2=1.0\,\mathrm V\)，
内阻 \(r_1=0.5\,\Omega\)、
\(r_2=1.0\,\Omega\)，
及 \(R_1=4.5\,\Omega,R_2=19.0\,\Omega,
 R_3=10.0\,\Omega,R_4=5.0\,\Omega\)。
依据图16.22中左右结点、两电源极性、
中间结点与各电阻连接，计算全电路电流分布。

\paragraph{习题 16.13}
对图16.23的惠斯通电桥，
用四个外臂阻值 \(R_1,R_2,R_3,R_4\)
及检流计支路电流 \(I_g\) 写出三个回路
基尔霍夫方程，消去有关电流，
证明当
\[
 \frac{R_1}{R_2}=\frac{R_3}{R_4}
\]
时 \(I_g=0\)。

\paragraph{习题 16.14}
按图16.24晶体管电路，
电源 \(\mathcal E=6\,\mathrm V\)，
电源内阻零；
\(U_{ec}=1.96\,\mathrm V\)，
\(U_{eb}=0.2\,\mathrm V\)，
\(I_c=2\,\mathrm{mA}\)，
\(I_b=20\,\mu\mathrm A\)，
\(I_2=0.4\,\mathrm{mA}\)，
\(R_c=1\,\mathrm{k}\Omega\)。
求 \(R_1,R_2,R_e\)，
并依原图给出支路电流方向。

\paragraph{习题 16.15}
证明由电容 \(C\) 经电阻 \(R\)
放电时电阻耗散的总能量等于放电前
电容器储能。
回应「只有在 \(t=\infty\) 才完全放电，
所以不可能真正放完」的疑问：
给定合理 \(R,C,U_0\) 时，计算电荷降至
剩下一个电子电量的时间。

\paragraph{习题 16.16}
红宝石激光器的脉冲闪光灯使用
\(2000\,\mu\mathrm F\) 电容器，充至
\(4000\,\mathrm V\) 后放电。
若充电电源最大输出电流为 \(1\,\mathrm A\)，
求充电电路最小时间常数；
闪光灯管内电阻近似
\(0.5\,\Omega\)，
求最大放电电流及放电电路时间常数。

\paragraph{习题 16.17}
大型电磁铁在电压 \(400\,\mathrm V\)
下流过 \(200\,\mathrm A\) 电流，
用水冷却线圈。
若冷却水入口为 \(20^\circ\mathrm C\)，
出口不得超过 \(80^\circ\mathrm C\)，
求最小水流量，单位
\(\mathrm{L/min}\)。
忽略向环境散热，水的密度约
\(1\,\mathrm{kg/L}\)，
比热约 \(4.18\,\mathrm{kJ/(kg\,K)}\)。

\paragraph{习题 16.18}
利用恒定电流下电荷守恒、均匀导体
欧姆本构关系及高斯定律证明：
\emph{均匀导体内的稳恒电流不允许体自由电荷
在内部积聚}；
净自由电荷只能出现在导体表面
或不同导体/材料的界面处。
